\documentclass[10pt,a4paper]{amsart}
\usepackage{amsmath,amssymb,amsthm,graphicx}
\usepackage{verbatim}
\newenvironment{reference}{\comment}{\endcomment}
%\newenvironment{reference}{}{}
\newenvironment{slogan}{\comment}{\endcomment}
\newenvironment{history}{\comment}{\endcomment}

% For commutative diagrams we use Xy-pic
\usepackage[all]{xy}

% We use 2cell for 2-commutative diagrams.
\xyoption{2cell}
% Theorem environments.
%
\theoremstyle{plain}
\newtheorem{theorem}[subsection]{Theorem}
\newtheorem{proposition}[subsection]{Proposition}
\newtheorem{lemma}[subsection]{Lemma}

\theoremstyle{definition}
\newtheorem{definition}[subsection]{Definition}
\newtheorem{example}[subsection]{Example}
\newtheorem{exercise}[subsection]{Exercise}
\newtheorem{situation}[subsection]{Situation}

\theoremstyle{remark}
\newtheorem{remark}[subsection]{Remark}
\newtheorem{remarks}[subsection]{Remarks}

\numberwithin{equation}{subsection}

% Macros
%
\def\lim{\mathop{\mathrm{lim}}\nolimits}
\def\colim{\mathop{\mathrm{colim}}\nolimits}
\def\Spec{\mathop{\mathrm{Spec}}}
\def\Hom{\mathop{\mathrm{Hom}}\nolimits}
\def\Ext{\mathop{\mathrm{Ext}}\nolimits}
\def\SheafHom{\mathop{\mathcal{H}\!\mathit{om}}\nolimits}
\def\SheafExt{\mathop{\mathcal{E}\!\mathit{xt}}\nolimits}
\def\Sch{\mathit{Sch}}
\def\Mor{\mathop{\mathrm{Mor}}\nolimits}
\def\Ob{\mathop{\mathrm{Ob}}\nolimits}
\def\Sh{\mathop{\mathit{Sh}}\nolimits}
\def\NL{\mathop{N\!L}\nolimits}
\def\CH{\mathop{\mathrm{CH}}\nolimits}
\def\proetale{{pro\text{-}\acute{e}tale}}
\def\etale{{\acute{e}tale}}
\def\QCoh{\mathit{QCoh}}
\def\Ker{\mathop{\mathrm{Ker}}}
\def\Im{\mathop{\mathrm{Im}}}
\def\Coker{\mathop{\mathrm{Coker}}}
\def\Coim{\mathop{\mathrm{Coim}}}

% Boxtimes
%
\DeclareMathSymbol{\boxtimes}{\mathbin}{AMSa}{"02}

%
% Macros for moduli stacks/spaces
%
\def\QCohstack{\mathcal{QC}\!\mathit{oh}}
\def\Cohstack{\mathcal{C}\!\mathit{oh}}
\def\Spacesstack{\mathcal{S}\!\mathit{paces}}
\def\Quotfunctor{\mathrm{Quot}}
\def\Hilbfunctor{\mathrm{Hilb}}
\def\Curvesstack{\mathcal{C}\!\mathit{urves}}
\def\Polarizedstack{\mathcal{P}\!\mathit{olarized}}
\def\Complexesstack{\mathcal{C}\!\mathit{omplexes}}
% \Pic is the operator that assigns to X its picard group, usage \Pic(X)
% \Picardstack_{X/B} denotes the Picard stack of X over B
% \Picardfunctor_{X/B} denotes the Picard functor of X over B
\def\Pic{\mathop{\mathrm{Pic}}\nolimits}
\def\Picardstack{\mathcal{P}\!\mathit{ic}}
\def\Picardfunctor{\mathrm{Pic}}
\def\Deformationcategory{\mathcal{D}\!\mathit{ef}}
\usepackage{cleveref}
\crefname{originallemma}{original lemma}{original lemmas}
\Crefname{originallemma}{Original lemma}{Original lemmas}
\makeatletter
\newlabel{lemma-behaviour-of-dimensions-wrt-smooth-morphisms}{{0DRI}{source}{}{sourcepointer-1}{}}
\newlabel{lemma-behaviour-of-dimensions-wrt-smooth-morphisms@cref}{{[originallemma][0][]0DRI}{[1][0][]source}{}{}{}}
\makeatother
\crefname{originallemma}{original lemma}{original lemmas}
\Crefname{originallemma}{Original lemma}{Original lemmas}
\makeatletter
\newlabel{lemma-relative-dimension-is-semi-continuous}{{0DRQ}{source}{}{sourcepointer-2}{}}
\newlabel{lemma-relative-dimension-is-semi-continuous@cref}{{[originallemma][0][]0DRQ}{[1][0][]source}{}{}{}}
\makeatother
\crefname{originallemma}{original lemma}{original lemmas}
\Crefname{originallemma}{Original lemma}{Original lemmas}
\makeatletter
\newlabel{lemma-constancy-of-dimension}{{0DRT}{source}{}{sourcepointer-3}{}}
\newlabel{lemma-constancy-of-dimension@cref}{{[originallemma][0][]0DRT}{[1][0][]source}{}{}{}}
\makeatother
\title{Dimension is constant on an irreducible Jacobson pseudo-catenary locally Noetherian stack}
\author{The Stacks Project Authors}
\date{stacks-project:stacks-geometry.tex:local-snapshot-20261012; GNU FDL 1.2 or later}
\begin{document}
\maketitle
\noindent\textit{Editorial note.} All delivered mathematical text is the author's own, copied verbatim from the pinned Stacks Project chapter file stacks-geometry.tex; no statement or proof step was rewritten, reordered or removed.\par\medskip
\noindent\textit{Editorial note.} Span recomputation by label pairing (LANE\_BRIEF 10.2): in this chapter the author's proof environment follows the statement environment as a standalone block; the spec covers the complete environment pair (statement span, proof span) as recorded in the pool candidate.\par\medskip
\noindent\textit{Editorial note.} Occupation scan: combined claim\_key + statement\_anchor + generic-line scan over all registration files (active and withdrawn) found 0 hits for stacks-geometry.tex (ge-occupied.json, 2026-10-14); 207 route=geometry / chapter-list field hits in focused-seeds are alias noise (route name and alphabetical chapter enumerations), excluded per the ladder alias-blind-spot caveat.\par\medskip
\ifdefined\hypertarget\hypertarget{sourcepointer-1}{}\fi
\noindent\textit{Source reference.} Original lemma 0DRI: sources/stacks-chow/tags.txt, source lines 17223--17223. Cross-chapter standard prerequisite cited by the author inside the counted span; it is used as a pointer and its own proof is not part of this item.\par\medskip
\ifdefined\hypertarget\hypertarget{sourcepointer-2}{}\fi
\noindent\textit{Source reference.} Original lemma 0DRQ: sources/stacks-chow/tags.txt, source lines 17230--17230. Cross-chapter standard prerequisite cited by the author inside the counted span; it is used as a pointer and its own proof is not part of this item.\par\medskip
\ifdefined\hypertarget\hypertarget{sourcepointer-3}{}\fi
\noindent\textit{Source reference.} Original lemma 0DRT: sources/stacks-chow/tags.txt, source lines 17233--17233. Cross-chapter standard prerequisite cited by the author inside the counted span; it is used as a pointer and its own proof is not part of this item.\par\medskip
\begin{lemma}
\label{lemma-irreducible-implies-equidimensional}
If $\mathcal{X}$ is
a Jacobson, pseudo-catenary, and locally Noetherian  algebraic stack
for which $|\mathcal{X}|$ is irreducible,
then $\dim_x(\mathcal{X})$ is a constant function on $|\mathcal{X}|$.
\end{lemma}

\begin{proof}
It suffices to show that $\dim_x(\mathcal{X})$ is locally constant on
$|\mathcal{X}|$,
since it will then necessarily be constant (as $|\mathcal{X}|$ is connected,
being irreducible).  Since $\mathcal{X}$ is pseudo-catenary,
we may find a smooth surjective morphism $U \to \mathcal{X}$ with $U$
being a universally catenary scheme.  If $\{U_i\}$ is an
cover of $U$ by quasi-compact open subschemes, we may replace
$U$ by $\coprod U_i,$, and
it suffices to show that
the function $u \mapsto \dim_{f(u)}(\mathcal{X})$ is locally constant on $U_i$.
Since we check this for one $U_i$ at a time, we now drop the subscript,
and write simply $U$ rather than $U_i$.
Since $U$ is quasi-compact, it
is the union of a finite number of irreducible components,
say $T_1 \cup \ldots \cup T_n$.  Note that each $T_i$ is Jacobson,
catenary, and locally Noetherian,
being a closed subscheme of the Jacobson, catenary, and locally Noetherian
scheme $U$.

\medskip\noindent
By Lemma \ref{lemma-behaviour-of-dimensions-wrt-smooth-morphisms}, we have
$\dim_{f(u)}(\mathcal{X}) = \dim_{u}(U) - \dim_{u}(U_{f(u)}).$
Lemma \ref{lemma-relative-dimension-is-semi-continuous} (2)
shows that the second term in the right hand expression is locally
constant on $U$, as $f$ is smooth,
and hence we must show that $\dim_u(U)$
is locally constant on $U$.  Since $\dim_u(U)$ is the maximum
of the dimensions $\dim_u T_i$, as $T_i$ ranges over the components
of $U$ containing $u$, it suffices to show
that if a point $u$ lies on two distinct components,
say $T_i$ and $T_j$ (with $i \neq j$),
then $\dim_u T_i = \dim_u T_j$,
and then to note that $t\mapsto \dim_t T$ is a constant
function on an irreducible Jacobson,
catenary, and locally Noetherian scheme $T$
(as follows from Lemma \ref{lemma-constancy-of-dimension}).

\medskip\noindent
Let $V = T_i \setminus (\bigcup_{i' \neq i} T_{i'})$
and $W = T_j \setminus (\bigcup_{i' \neq j} T_{i'})$.
Then each of $V$ and $W$ is a non-empty open subset of $U$,
and so each has non-empty open image in $|\mathcal{X}|$.  As $|\mathcal{X}|$ is
irreducible,
these two non-empty open subsets of $|\mathcal{X}|$ have a non-empty
intersection.
Let $x$ be a point lying in this intersection, and let $v \in V$ and
$w\in W$ be points mapping to $x$.
We then find that
$$
\dim T_i = \dim V = \dim_v (U) = \dim_x (\mathcal{X}) + \dim_v (U_x)
$$
and similarly that
$$
\dim T_j = \dim W = \dim_w (U) = \dim_x (\mathcal{X}) + \dim_w (U_x).
$$
Since $u \mapsto \dim_u (U_{f(u)})$ is locally constant on $U$,
and since $T_i \cup T_j$ is connected (being the union of two irreducible,
hence connected, sets that have non-empty intersection),
we see that $\dim_v (U_x) = \dim_w(U_x)$,
and hence, comparing the preceding two equations,
that $\dim T_i = \dim T_j$, as required.
\end{proof}

\end{document}
